\section{Future work}
\label{sec:future_work}
\label{sec:research-plan}

% Davide: one or two sentences on where the project stands at the
% assessment date go here, once the benchmark and the WASP-17 b analysis
% are closed. The rest of the section does not depend on that wording.

In total, I plan to publish six papers that all tie back to NOVA. The
first presents NOVA itself, the injection benchmark and the WASP-17~b
analysis, including a comparison with published transmission spectra. The
next two apply NOVA to other selected SOSS targets, addressing transfer to
other targets, repeated visits of the same planet and stellar variability.
Once that is done, I aim to reanalyse a broad range of published SOSS
transmission spectra with NOVA. The final two papers extend NOVA to
NIRSpec and to secondary eclipses. Figure~\ref{fig:gantt} shows the
planned schedule.

\subsection{Paper 1: NOVA, the benchmark and WASP-17~b}
\label{subsec:fw_benchmark}

\noindent\textbf{Research question:} \textit{Does NOVA recover
transmission spectra more accurately, and with more reliable
uncertainties, than extraction-based pipelines, and does it change what we
conclude about the atmosphere of WASP-17~b?}

% TODO (Davide): one paragraph on what remains for Paper 1 after the
% assessment (submission quarter, benchmark package release, whether the
% six-visit matrix is in). Depends on what is finished in the next month.

\subsection{Paper 2: Transfer to other SOSS targets}
\label{subsec:fw_transfer}

\noindent\textbf{Research question:} \textit{Do the atmospheric
conclusions for other SOSS targets change when their observations are
analysed with NOVA, and if so, why?}

I will start with WASP-39~b and WASP-96~b. Both are hot gas giants with
well-studied SOSS observations, so the first transfer stays close to
WASP-17~b. WASP-39~b has strong water features and published constraints
on its water abundance and clouds \citep{FeinsteinEtAl2023,CarterEtAl2024}.
WASP-96~b is interesting because the published analyses disagree:
\citet{TaylorEtAl2023} found an enhanced scattering slope at short
wavelengths, whereas \citet{WangEtAl2026} found no sign of haze. Both
fields contain contaminating stars.

% ============================================================
% TODO (Davide): BIG ONE. The three contamination paragraphs below
% describe the state on 14 September 2026 and the contamination policy
% is under active work with Astra. Re-verify every factual statement
% before the ESA is handed in, in particular:
%   - whether contaminant pixels inside the aperture are excluded from
%     the fit by default (three-arm aperture test pending);
%   - whether the injector's source calibration excludes contaminant
%     light (Full-light / Penumbra work, in progress);
%   - the fairness sentence (same injected data, own native processing);
%   - the WASP-96 b template prototype claim.
% ============================================================
\paragraph{Field-star contamination.}
NOVA has no explicit model for contaminating stars. Two parts of the
method are insensitive to their light. The spatial response $q^\star$ is
estimated from how much each pixel dims during transit
(Equation~\ref{eq:qstar-fit}), so the steady light of a field star ends up
in the constant term of that fit and never enters $q^\star$. The
background is a smooth surface fitted on pixels far from the traces and
from any detected contaminant, so it cannot bend into the shape of a
contaminant trace. The transit fit itself is not protected. Contaminant
light that lands under the target trace stays in the data, the only model
term that can absorb it is the star, and the measured transit depth comes
out too shallow. \citet{LouieEtAl2025} show for WASP-17~b how hard a
correction is when only a small part of the contaminant trace is visible;
JExoRES shows that a template correction works when more of it is visible
\citep{HolmbergMadhusudhan2023}.

NOVA applies no field-star mask to the pixels it fits: a pixel is
excluded only by the fitting aperture, the detector validity flags and
the calibrated wavelength support. The comparison with other pipelines is
fair not because every pipeline fits the same pixels, but because all of
them receive the same injected detector data and then apply their own
declared processing, including their own contaminant handling.
Contaminants are located in two ways: from the science frames themselves,
by segmenting the median image, and from the F277W exposure, in which the
filter blocks the target spectrum below about 2.4~$\mu$m so that field
stars and parts of their traces are seen directly. Their brightness
cannot simply be transferred from the F277W image to the science
exposures, because the filter changes the transmitted light
\citep{HolmbergMadhusudhan2023}. Such maps exist for all six benchmark
visits.

These contaminant maps already enter NOVA in two places that do not touch
the fitted pixels. The off-trace pixels on which the background is
calibrated exclude the detected contaminants with a margin around each,
together with the trace cores and their wings. The $1/f$ noise estimate
rejects the F277W contaminant pixels, as exoTEDRF does natively. The
third place is the injector, and this one is still open: the stellar
source model must be calibrated so that contaminant light is not injected
as target light. For the fit itself, where a contaminant has to be
treated, I will first simply exclude the affected pixels, or whole columns
where the map is not precise enough, and report the lost wavelength
coverage. Where that loss matters for the science, I will add an additive
template per contaminant to the model. Its shape follows the visible part
of the contaminant trace, projected under the target with the instrument
model; its amplitude is a nuisance parameter anchored like the background,
so it cannot trade against the transit depth. A first prototype of such a
template exists for WASP-96~b.

\paragraph{Transfer.}
NOVA itself does not change: the transit model, the detector likelihood,
the estimation of $q^\star$ and the covariance carry over as they are.
What changes from visit to visit is the input: transit geometry, limb
darkening, integration timing, excluded pixels, trace position and
background. I will estimate these from each new observation with the same
procedures as for WASP-17~b. The benchmark injections of Paper 1 already
test whether those procedures recover known spectra in these two detector
scenes, including weak features. If a procedure has to change, I will
record the change and report what it does to the WASP-17~b result.

\paragraph{Testing the transfer.}
I will test the transfer with injection tests of the same kind as in
Paper 1. The procedure is fixed on WASP-39~b and then applied unchanged
to WASP-96~b, as a test of transfer to a different detector scene.

\paragraph{Atmospheric interpretation.}
For each planet I will compare the NOVA spectrum with the published
spectra and, where those pipelines are available, with my own runs of the
pipelines used in the published analyses on the same data. The same
pipelines are also run on the injection tests for that target, so that
their accuracy in that detector scene is known before the real spectra
are compared. The same atmospheric retrieval is then run on all spectra,
to check whether the water abundance, the cloud constraints and the
spectral slopes agree. WASP-96~b is the clearest case: its two published
analyses used different pipelines and reached different conclusions about
the slope, so this comparison tests directly whether the disagreement
comes from the reduction. The outcome may support either published
interpretation, or show that the slope cannot be decided from these data.

\subsection{Paper 3: Repeated visits and stellar variability}
\label{subsec:fw_repeats}

\noindent\textbf{Research question:} \textit{Do repeated visits of the
same planet give the same spectrum once instrumental and stellar effects
are accounted for, and if not, do the error bars explain the difference?}

I will start with GJ 9827~d. It has two SOSS visits and a published
water-rich interpretation \citep{PiauletGhorayebEtAl2024}, it is a much
smaller planet than the gas giants of Paper 2, and two visits give a
direct test of repeatability. I will first analyse each visit on its own,
excluding the integrations affected by flares, and compare the two
spectra. I will then let the visits share the orbital geometry and see how
much that improves their agreement. I will also run injection tests on
both visits, since agreement between two visits does not by itself show
that either spectrum is right.

Next I will add a flare model to the fit instead of throwing the flared
integrations away. The flare happens at the same moment at every
wavelength, so its timing is shared across the spectrum, but its
brightness is fitted separately at each wavelength. The point is to find
out whether the integrations that masking discards still contain useful
information about the planet. One complication is specific to NOVA:
$q^\star$ is estimated from how each pixel dims during transit, and a
flare changes the brightness in the same time series, so it can distort
$q^\star$ before the spectrum is even fitted. I will therefore redo the
$q^\star$ estimation with the flare model included, and test other ways of
estimating $q^\star$ for cases where the transit is too weak to pin it
down on its own. If the flare model works, it becomes a standard part of
NOVA and is available for every later target.

The second main target is LHS 1140~b, which has four SOSS visits spread
over different dates. Flares are not the only stellar effect. Dark spots
and bright faculae change the apparent transmission spectrum, because the
part of the star the planet crosses is not the same as the average of the
visible disc \citep{RackhamEtAl2018,RackhamEtAl2019}, and a planet
crossing a spot leaves a bump in the light curve. Where GJ 9827~d tests
the flare model, LHS 1140~b is the case for the rest of the stellar
model: unocculted spots and faculae, and an activity level that changes
from visit to visit, modelled separately from flares and from
instrumental trends. Together the two targets are meant to leave NOVA
with a stellar-activity component that later targets can use without new
development. The two most recent analyses agree that helium is
not detected \citep{GressierEtAl2026,BennettEtAl2026}, while the first
study reported tentative absorption features, favouring a nitrogen-rich
atmosphere at low significance, together with signs of stellar
contamination \citep{CadieuxEtAl2024}. I will check whether those
tentative features are still there once the stellar changes between
visits are accounted for. Where a visit has a known complication, such as
a displaced trace or a second planet transiting during the observation, I
will model it explicitly.

TRAPPIST-1~f is the hardest case and the stress test for the combined
model. It has five SOSS visits, every one of them with flares, although
the published spectrum shows no sign of contamination from unocculted
spots \citep{LimEtAl2026}. The question is whether the stellar model
built on the first two systems copes with these data while leaving the
planetary signal intact. Some visits may have too little clean data and
give only upper limits or a reduced wavelength range.

\subsection{Paper 4: A uniform reanalysis of SOSS transmission spectra}
\label{subsec:fw_atlas}

\noindent\textbf{Research question:} \textit{Once NOVA is robust, which
atmospheric features change when a broad range of published SOSS
transmission spectra is reanalysed with it, and which published
conclusions hold up?}

Once Papers 2 and 3 have shown that NOVA transfers between targets and
copes with stellar activity, I will apply it to a broad, selected sample
of the public SOSS transmission observations. Uniform reanalyses of this
kind have been done with Hubble data \citep{TsiarasEtAl2018,SabaEtAl2025}.
When every planet is analysed with the same pipeline, a difference between
two spectra is a difference between two planets and not between two
reductions, which makes trends across the population visible. This
catalogue does the same for SOSS.

The sample deliberately includes cases that NOVA was not developed on.
HAT-P-18~b and WASP-52~b have active hosts, with unocculted spots in one
case and spot crossings in the other
\citep{FournierTondreauEtAl2024,FournierTondreauEtAl2025}. GJ 357~b and
LTT 9779~b have small or muted signals
\citep{TaylorEtAl2025,RadicaEtAl2024}. K2-18~b has already been studied
for how much its spectrum depends on reduction and modelling choices, so
it gives a direct comparison \citep{SchmidtEtAl2025}. HAT-P-26~b and
WASP-80~b are two of the six benchmark hosts, so their detector scenes
are already characterised, and together with HAT-P-18~b they are the
first entries of the catalogue.

A catalogue of this size needs a workflow that takes an observation from
the raw files to a validated spectrum without anyone editing code along
the way: locating the contaminating stars, fitting the spectrum, running
the injection tests. The contamination treatment is the one settled in
Paper 2, applied to each observation with its own source positions and
masks. For a subset of targets, and for every case where NOVA changes an
important conclusion, I will also run a conventional reduction of the
same data for comparison.

Every spectrum in the catalogue comes with its own injection tests, of
the same kind as in Paper 1. The tests may show that a spectrum is
reliable at most wavelengths but not at all of them, for example where a
contaminant crosses the trace. Those wavelengths are flagged in the
released spectrum.

For each planet I will compare the NOVA spectrum with the published one
at two levels: first on the spectra themselves, asking whether the
absorption features or the slope of the spectrum change, and then with
retrievals, asking whether the inferred gas abundances and cloud
properties change.

\subsection{Paper 5: Joint NIRISS and NIRSpec analysis}
\label{subsec:fw_nirspec}

\noindent\textbf{Research question:} \textit{Once NOVA is extended to
NIRSpec, can observations of the same planet with both instruments be
explained by one transmission spectrum?}

A first test will be WASP-39~b. It was observed with four JWST instrument
modes, and \citet{CarterEtAl2024} reanalysed all four with the same
orbital parameters, limb darkening and treatment of systematics, after
which the spectra agreed much better where their wavelengths overlap.
Better agreement is not the same as a correct spectrum, though: once
every data set is analysed with the same choices, the spectra will agree
whether or not those choices are right. So I will build injection tests
for the NIRSpec observations of WASP-39~b, as Paper 1 does for SOSS,
extend NOVA to the two NIRSpec modes, and check on injected spectra that
each extension recovers the truth. Only then will I compare the three
real spectra, SOSS and the two NIRSpec modes, and ask whether they agree.
If they do, the agreement means something, because each spectrum has
been validated on its own.
% TODO (Davide): the next paragraph still orders the NIRSpec modes as
% G235H on LHS 1140 b first, then G395H, then PRISM. WASP-39 b's NIRSpec
% modes are G395H and PRISM, so the order and the first target must change
% to match this paragraph. Not yet reviewed in the paragraph pass.

NIRSpec G235H/F170LP comes first because its band overlaps SOSS
substantially. LHS 1140~b already has G235H and G395H observations
\citep{DamianoEtAl2024} and is the initial target. G395H/F290LP then
extends the wavelength coverage, followed by dedicated PRISM tests and
selected G140 modes. The actual wavelength range and detector gap of each
configuration will be used; G395H mainly extends coverage and offers
little overlap with SOSS \citep{JDoxNIRSpecBOTS2026}. WASP-39~b PRISM is a
separate test of the documented disagreement.

The extension requires a NIRSpec description of how stellar light falls
on the detector: trace shape, wavelength mapping, spectral resolution,
slit throughput and noise, while the transit model is retained. The
injector uses the same instrument description to create tests with known
spectra, and each mode is validated separately before it enters a joint
analysis.

NIRISS and NIRSpec are first fitted independently and compared at matched
resolution in the overlap. The joint model then uses one intrinsic
spectrum with separate backgrounds, detector effects and stellar
properties for each visit. Stellar changes between dates can shift the
apparent depths even when the atmosphere is unchanged, so those
contributions are constrained by physical models and by the data rather
than by an arbitrary offset at every wavelength, which could conceal a
disagreement. The common spectrum is passed through each instrument's
response, so that the comparison accounts for the different bins and
resolutions.

Paired injections test recovery of the same planet through both
instruments with independent noise, backgrounds and stellar states, and
deliberate mismatches test whether the joint model detects incompatible
data. The products include individual and joint spectra with the
uncertainty of relative depth offsets, and data sets that remain
inconsistent are reported as such. The first NIRSpec paper covers the
modes that pass these tests; further modes follow in later releases.

\subsection{Paper 6: Secondary eclipses and the dayside atmosphere}
\label{subsec:fw_eclipses}

\noindent\textbf{Research question:} \textit{Can NOVA improve the
reliability of weak eclipse measurements and the atmospheric conclusions
drawn from them?}

WASP-17~b is the starting point. Its SOSS eclipse analysis found that
allowing negative fitted depths changed a tentative FeH constraint while
the water result was more stable \citep{GressierEtAl2024}. I will
investigate whether a detector-level analysis improves these weak
measurements and how the resulting spectrum changes the estimated dayside
temperature and composition.

During secondary eclipse the planet's thermal emission and reflected light
disappear while the stellar light remains. The planet will initially be a
uniformly bright disc whose hidden fraction is predicted through ingress
and egress and averaged over each integration. The detector model keeps
its treatment of source light and background, with an eclipse model
replacing the transit signal. Uncertainty in the orbit and eclipse timing,
uneven planetary brightness and slow changes of planet flux around the
event will be tested.

The positive depth parameters become signed eclipse amplitudes. Noise can
place a weak measurement below zero; retaining that estimate lets the
atmospheric fit use the full measurement uncertainty, while the physical
atmosphere still predicts nonnegative flux. Depth relative to the star is
distinguished from depth relative to star and planet together. Estimating
the spatial response $q^\star$ needs its own tests, since a weak eclipse
may contain too little variation to calibrate it through the event; the
mean stellar image and the instrument's optical model are the
alternatives.

The eclipse injector adds a known planetary contribution to a star-only
baseline. The added light disappears during eclipse, and star, background
and unrelated sources are unchanged. For injection before calibration,
the accumulated planetary counts are added to each detector read with the
corresponding photon noise, which changes the signal and noise treatment
of the transit injector. Tests cover very weak emission, thermal and
reflected light, and atmospheres with and without weak molecular
features. Conventional extraction and NOVA are compared with signed and
positive-constrained fits, followed by matched retrievals, to separate
the effect of the depth constraint from the effect of the reduction. The
metrics are recovery of injected spectra, coverage of uncertainty
intervals and false detections of molecular features. The validated
procedure is then applied to WASP-17~b. Circulation models can support
later phase-curve tests.

\subsection{Timeline, dependencies and risks}
\label{subsec:fw_timeline}

\begin{figure*}[t]
\centering
\resizebox{\textwidth}{!}{%
\begin{ganttchart}[
  hgrid,
  vgrid,
  x unit=0.84cm,
  y unit title=0.55cm,
  y unit chart=0.52cm,
  title/.style={fill=black!6,draw=black!25},
  title label font=\bfseries\scriptsize,
  bar label font=\scriptsize,
  group label font=\bfseries\scriptsize,
  milestone label font=\scriptsize,
  bar height=0.55,
  group height=0.34,
  group peaks height=0.16,
  milestone height=0.45
]{1}{13}
\gantttitle{2026}{1}\gantttitle{2027}{4}\gantttitle{2028}{4}\gantttitle{2029}{4}\\
\gantttitle{Q4}{1}
\gantttitle{Q1}{1}\gantttitle{Q2}{1}\gantttitle{Q3}{1}\gantttitle{Q4}{1}
\gantttitle{Q1}{1}\gantttitle{Q2}{1}\gantttitle{Q3}{1}\gantttitle{Q4}{1}
\gantttitle{Q1}{1}\gantttitle{Q2}{1}\gantttitle{Q3}{1}\gantttitle{Q4}{1}\\
\ganttgroup[group/.append style={fill=ganttsoss!35,draw=ganttsoss}]
  {SOSS transmission programme}{1}{8}\\
\ganttbar[bar/.append style={fill=ganttsoss!70,draw=ganttsoss}]
  {Paper 1: submission and release}{1}{2}\\
\ganttmilestone[milestone/.append style={fill=ganttsoss,draw=ganttsoss}]
  {Paper 1 submitted}{1}\\
\ganttbar[bar/.append style={fill=ganttsoss!70,draw=ganttsoss}]
  {Paper 2: WASP-39 b, WASP-96 b}{2}{4}\\
\ganttbar[bar/.append style={fill=ganttsoss!70,draw=ganttsoss}]
  {Paper 3: repeated visits, variability}{4}{6}\\
\ganttbar[bar/.append style={fill=ganttsoss!70,draw=ganttsoss}]
  {Paper 4: uniform SOSS reanalysis}{6}{8}\\
\ganttmilestone[milestone/.append style={fill=ganttsoss,draw=ganttsoss}]
  {SOSS catalogue release}{8}\\
\ganttgroup[group/.append style={fill=ganttgeneral!35,draw=ganttgeneral}]
  {Instrument and geometry extensions}{7}{10}\\
\ganttbar[bar/.append style={fill=ganttgeneral!70,draw=ganttgeneral}]
  {Paper 5: NIRISS and NIRSpec}{7}{9}\\
\ganttbar[bar/.append style={fill=ganttgeneral!70,draw=ganttgeneral}]
  {Paper 6: secondary eclipses}{9}{10}\\
\ganttmilestone[milestone/.append style={fill=ganttgeneral,draw=ganttgeneral}]
  {Science freeze}{10}\\
\ganttgroup[group/.append style={fill=ganttwrite!35,draw=ganttwrite}]
  {Thesis}{2}{13}\\
\ganttbar[bar/.append style={fill=ganttwrite!70,draw=ganttwrite}]
  {Thesis integration}{2}{10}\\
\ganttbar[bar/.append style={fill=ganttwrite!70,draw=ganttwrite}]
  {Formal writing-up period}{11}{13}\\
\ganttmilestone[milestone/.append style={fill=ganttwrite,draw=ganttwrite}]
  {Thesis submission}{13}
\end{ganttchart}%
}
\caption{Publication and thesis plan in quarters, starting after the
assessment. Paper 1 is submitted in 2026 Q4. The SOSS transmission
programme closes with the catalogue release in 2028 Q3, and the instrument
and geometry extensions end with a science freeze in 2029 Q1. The research
phase ends on approximately 27 April 2029 after the one-month
interruption, and the formal writing-up period runs to 28 October 2029.}
\label{fig:gantt}
\end{figure*}

Figure~\ref{fig:gantt} shows the dependency structure. Paper 1 releases
the benchmark and the accepted injector that every later paper reuses.
Paper 2 fixes the contamination policy that the catalogue applies without
change, and Paper 3 supplies the stellar-variability model that Papers 4
and 5 need for spotted hosts and multi-epoch data. The NIRSpec detector
description of Paper 5 is new development and starts while the catalogue
is being produced. Paper 6 reuses the injector changes of Paper 5 for
injection before calibration. In every paper the recovered spectra are
frozen before the atmospheric retrievals are run, so that detector-model
uncertainty is traced into the atmospheric constraints rather than fitted
away. The code and validation products of each paper are archived with a
citable identifier. Further extensions to MIRI, phase curves, direct
atmospheric retrieval from detector data and accelerated inference depend
on these results and are not scheduled. Progress is reviewed quarterly
against closed deliverables, and a milestone counts as complete only when
its data products, configuration, diagnostics and written interpretation
are archived together.

Five risks are foreseen. The first is contamination: if excluded pixels
remove more wavelength coverage than a scientific question can tolerate
and the template correction does not validate, the affected spectra are
released with the reduced coverage and the limitation is stated. The
second is stellar variability: where flares or spot crossings cannot be
separated from the planetary signal, the analysis reports limits or a
restricted wavelength range rather than a spectrum. The third is
computational cost, which grows with the number of visits and injection
tests; variable projection, streamed derivatives and staged small cases
before visit-scale runs control it, as in Paper 1. The fourth is the
NIRSpec detector description, the largest technical development of the
later papers; if it does not validate in time, Paper 5 is reduced to
independent fits compared in the overlap region. The fifth is schedule: a
delay in the referee process or the benchmark release of Paper 1 shifts
Paper 2, and Papers 5 and 6 are the adjustable scope, while the SOSS
programme of Papers 1 to 4 is the core of the thesis. A fair result in
which NOVA does not outperform extraction-based analyses remains
scientifically useful, because it bounds the benefit of joint
detector-level inference under stated conditions.
